% Source: https://artofproblemsolving.com/community/c3h2317350_100th_post
% Stable identity: aops-topic-2317350-q01
The angle bisectors $AF$ and $CE$ meet at the incentre $I$, so $I$ lies on $BD$. The bisector of $\angle ABC$ also passes through $B$ and $I$; since $I\ne B$, it is the line $BI=BD$. So the altitude $BD$ bisects $\angle ABC$.

Then the right triangles $ABD$ and $CBD$ have the common side $BD$, right angles at $D$ and equal angles at $B$, so they are congruent (ASA). Hence $BA=BC$ and $AD=DC$: the triangle is isosceles and $D$ is the midpoint of $AC$, so $BD$ is a median.

\begin{asy}
usepackage("amsmath");
size(5cm);
defaultpen(fontsize(10pt));
pair A=(0,0), C=(6,0), B=(3,9);
pair D=(A+C)/2, M=(B+C)/2, H=(A+B+C)/3;
real a=abs(B-C), b=abs(C-A), c=abs(A-B);
pair I=(a*A+b*B+c*C)/(a+b+c);
pair Fp=extension(A,I,B,C), Ep=extension(C,I,A,B);
draw(A--B--C--cycle,linewidth(0.9));
draw(B--D,blue); draw(A--Fp,dashed); draw(C--Ep,dashed); draw(A--M,red);
draw(D+(0,0.35)--D+(0.35,0.35)--D+(0.35,0));
dot(I); dot(H,red); dot(M); dot(Ep); dot(Fp);
label("$A$",A,SW); label("$C$",C,SE); label("$B$",B,N); label("$D$",D,S);
label("$E$",Ep,W); label("$F$",Fp,E); label("$M$",M,E); label("$I$",I+(0.8,0)); label("$H$",H,NW,red);
\end{asy}

The median $AM$ from $A$ (with $M$ the midpoint of $BC$) meets the median $BD$ at the centroid $H$, which divides each median in the ratio $2:1$ from the vertex. Therefore
\[\dfrac{BH}{HD}=\boxed{2}.\]
